\documentclass[11pt,a4paper]{article}
\usepackage{ctex}
\usepackage{amsmath,amssymb,bm}
\usepackage{verbatim}
\usepackage{array}
\usepackage{geometry}
\usepackage{float}
\usepackage{booktabs}
\usepackage{enumitem}
\usepackage{listings}
\usepackage{xcolor}

% 页面边距设置
\geometry{left=2.4cm,right=2.4cm,top=2.5cm,bottom=2.5cm}

% Verilog代码高亮配置
\lstset{
    language=Verilog,
    basicstyle=\ttfamily\small,
    commentstyle=\color{gray},
    keywordstyle=\color{blue},
    stringstyle=\color{red},
    frame=single,
    numbers=left,
    numberstyle=\tiny,
    breaklines=true,
    tabsize=4,
    escapeinside={/*@}{@*/},
    columns=fixed
}

% 枚举间距全局规范化
\setlist[enumerate,1]{itemsep=0.3em, topsep=0.2em}
\setlist[enumerate,2]{itemsep=0.2em, topsep=0.1em}
\setlist[itemize]{itemsep=0.2em}

\begin{document}
\title{实验三：基于FPGA实现组合逻辑电路实验指导书}
\author{}
\date{}
\maketitle

\section{实验内容总览}
本次实验包含\textbf{四项必做任务}与\textbf{两项选做拓展任务}，需严格按照流程依次完成，所有设计均基于 Xilinx EGO1（XC7A35T）FPGA 开发板。
\begin{enumerate}[label=任务\arabic*：]
    \item 基础示例：二输入与非门，拨码开关、按键作为输入，LED 作为输出；
    \item 必做：带使能端 2-4 译码器，拨码输入、四路 LED 输出；
    \item 必做：带使能端 3-8 译码器，拨码输入、八路 LED 输出；
    \item 必做：1 位全加器，三路拨码输入，和、进位分别由 LED 显示；
\end{enumerate}
\noindent \textbf{选做内容}：优先报警电路、4位减法运算+数码管显示电路。

\section{通用标准化实验流程（全任务通用）}
\subsection{阶段1：新建Vivado工程}
\begin{enumerate}
    \item 启动 Vivado 软件，点击 \textbf{Create New Project}，连续点击 \textbf{Next}；
    \item 工程名、存放路径\textbf{禁止包含中文、空格、特殊字符}，选择工程类型为 \textbf{RTL Project}；
    \item 器件选型：Artix-7 系列，封装 csg324，速度等级 -1，芯片型号 \verb|xc7a35t-1csg324-1|，点击 \textbf{Finish} 完成创建。
\end{enumerate}
\textbf{核查要点}：器件型号与开发板完全匹配。

\subsection{阶段2：添加设计源文件（Design Source）}
\begin{enumerate}
    \item 左侧 \texttt{Sources} 面板 $\rightarrow$ \texttt{Design Sources} $\rightarrow$ \texttt{Add Sources}；
    \item 选择 \textbf{Add or create design sources} $\rightarrow$ \textbf{Create File}，文件类型选择 Verilog；
    \item 自定义模块名与端口，粘贴对应功能代码，保存文件。
\end{enumerate}
\textbf{核查要点}：代码无语法报错，端口名称统一，与后续约束文件保持一致。

\subsection{阶段3：添加仿真源文件（Simulation Source）}
\begin{enumerate}
    \item 再次点击 \texttt{Add Sources}，选择 \textbf{Add or create simulation sources}；
    \item 新建仿真文件，命名规则：\texttt{模块名\_tb.v}，文件内例化待测模块并编写激励信号；
    \item 执行 \texttt{Simulation $\rightarrow$ Run Behavioral Simulation}，查看波形。
\end{enumerate}
\textbf{核查要点}：仿真波形与真值表完全一致，逻辑无误后方可进入硬件调试。

\subsection{阶段4：添加管脚约束文件（XDC）}
\begin{enumerate}
    \item \texttt{Add Sources} $\rightarrow$ \textbf{Add or create constraints}，新建约束文件 \texttt{const.xdc}；
    \item 写入三类约束：FPGA 管脚绑定、IO 电平标准 \texttt{LVCMOS33}、闲置引脚配置。
\end{enumerate}
\textbf{核查要点}：端口名、物理引脚编号对照 EGO1 板卡手册，无引脚冲突。

\subsection{阶段5：综合、布局布线、生成比特流}
\begin{enumerate}
    \item 执行 \textbf{Run Synthesis}（逻辑综合），保证无 Error 错误；
    \item 执行 \textbf{Run Implementation}（布局布线），保证无 Error 错误；
    \item 执行 \textbf{Generate Bitstream}，生成可下载的 \texttt{*.bit} 比特流文件。
\end{enumerate}
\textbf{核查要点}：三步流程全部正常完成，比特流文件生成成功。

\subsection{阶段6：JTAG在线下载与板级实测}
\begin{enumerate}
    \item 硬件设置：开发板跳线 \textbf{JP2} 接 \texttt{JTAG} 档位，Micro-USB 连接 J6 接口并上电；
    \item Vivado 中打开 \texttt{Hardware Manager} $\rightarrow$ \texttt{Open Target} $\rightarrow$ \texttt{Auto Connect}，自动识别 FPGA；
    \item 右键 FPGA 芯片 \texttt{xc7a35t\_0}，选择 \textbf{Program Device}，选中 \texttt{impl\_1} 下的 \texttt{.bit} 文件完成下载；
    \item 拨动拨码/按键，观察 LED 状态，对比理论逻辑。
\end{enumerate}
\textbf{核查要点}：硬件现象与真值表、仿真波形完全一致。

\subsection{可选拓展：程序固化至SPI Flash（断电保存）}
\begin{enumerate}
    \item 右键 \texttt{Generate Bitstream} $\rightarrow$ \texttt{Bitstream Setting}，勾选生成 \texttt{.bin} 文件，重新生成文件；
    \item 硬件管理器中选择 \textbf{Add Configuration Memory Device}，搜索并选择存储芯片 \texttt{n25q64-3.3V}；
    \item 右键 Flash 芯片，选择 \textbf{Program Configuration Memory Device}，烧写 \texttt{.bin} 文件；
    \item 断电后将 \textbf{JP2} 切至 \texttt{QSPI} 档位，重新上电，程序自动加载运行。
\end{enumerate}
\textbf{核查要点}：断电重启后电路功能保持不变。

\newpage
\section{分任务代码、引脚、约束与实验说明}
\subsection{任务1：二输入与非门（入门示例）}
\subsubsection{引脚分配}
输入：拨码开关 \textbf{SW7(P5)}、按键 \textbf{S1(R17)}；输出：LED \textbf{LD2\_7(F6)}。

\subsubsection{Verilog 设计代码}
\begin{lstlisting}
// 二输入与非门 数据流描述
module nand_test(
    input  SW7,
    input  S1,
    output LD2_7
);
assign LD2_7 = ~(SW7 & S1);
endmodule
\end{lstlisting}

\subsubsection{Verilog 仿真代码 \texttt{nand\_test\_tb.v}}
\begin{lstlisting}
module nand_test_tb();
reg SW7_tb, S1_tb;
wire LD2_7_tb;

// 例化待测模块
nand_test uut(
    .SW7(SW7_tb),
    .S1(S1_tb),
    .LD2_7(LD2_7_tb)
);

initial begin
    SW7_tb = 1'b0; S1_tb = 1'b0;
    #5 SW7_tb = 1'b0; S1_tb = 1'b1;
    #5 SW7_tb = 1'b1; S1_tb = 1'b0;
    #5 SW7_tb = 1'b1; S1_tb = 1'b1;
    #10 $finish;
end
endmodule
\end{lstlisting}

\subsubsection{XDC 约束文件}
\begin{lstlisting}
# 管脚绑定
set_property PACKAGE_PIN P5  [get_ports SW7]
set_property PACKAGE_PIN R17 [get_ports S1]
set_property PACKAGE_PIN F6  [get_ports LD2_7]
# 电平标准
set_property IOSTANDARD LVCMOS33 [get_ports *]
# 闲置引脚配置
set_property BITSTREAM.CONFIG.UNUSEDPIN PULLNONE [current_design]
\end{lstlisting}

\subsubsection{实验现象}
当 $\texttt{SW7}=1$ 且 $\texttt{S1}=1$ 时，LED $\texttt{LD2\_7}$ 熄灭；其余输入组合下 LED 点亮。

\newpage
\subsection{任务2：带使能端 2-4 译码器（必做）}
\subsubsection{引脚分配}
使能 \textbf{EN(M4)}、输入 $\boldsymbol{A_1}$(N4)、$\boldsymbol{A_0}$(R1)；
输出：$\boldsymbol{Y_0}$(K2)、$\boldsymbol{Y_1}$(J2)、$\boldsymbol{Y_2}$(J3)、$\boldsymbol{Y_3}$(H4)。

\subsubsection{真值表}
\begin{table}[htbp]
\centering
\begin{tabular}{ccccccc}
\toprule
$\text{EN}$ & $A_1$ & $A_0$ & $Y_3$ & $Y_2$ & $Y_1$ & $Y_0$ \\
\midrule
0  & $x$  & $x$  & 0  & 0  & 0  & 0  \\
1  & 0  & 0  & 0  & 0  & 0  & 1  \\
1  & 0  & 1  & 0  & 0  & 1  & 0  \\
1  & 1  & 0  & 0  & 1  & 0  & 0  \\
1  & 1  & 1  & 1  & 0  & 0  & 0  \\
\bottomrule
\end{tabular}
\end{table}

\subsubsection{Verilog 设计代码（行为级）}
\begin{lstlisting}
module Vr2to4dec_b(
    input        EN,
    input        in_A0,
    input        in_A1,
    output reg   out_Y0,
    output reg   out_Y1,
    output reg   out_Y2,
    output reg   out_Y3
);

always @(*) begin
    if(!EN) begin
        {out_Y3,out_Y2,out_Y1,out_Y0} = 4'b0000;
    end
    else begin
        case({in_A1, in_A0})
            2'b00: {out_Y3,out_Y2,out_Y1,out_Y0} = 4'b0001;
            2'b01: {out_Y3,out_Y2,out_Y1,out_Y0} = 4'b0010;
            2'b10: {out_Y3,out_Y2,out_Y1,out_Y0} = 4'b0100;
            2'b11: {out_Y3,out_Y2,out_Y1,out_Y0} = 4'b1000;
            default:{out_Y3,out_Y2,out_Y1,out_Y0} = 4'b0000;
        endcase
    end
end
endmodule
\end{lstlisting}

\subsubsection{Verilog 仿真代码 \texttt{Vr2to4dec\_tb.v}}
\begin{lstlisting}
module Vr2to4dec_tb();
reg ENt, A0t, A1t;
wire Y0t, Y1t, Y2t, Y3t;

Vr2to4dec_b uut(
    .EN(ENt),
    .in_A0(A0t),
    .in_A1(A1t),
    .out_Y0(Y0t),
    .out_Y1(Y1t),
    .out_Y2(Y2t),
    .out_Y3(Y3t)
);

initial begin
    {ENt,A1t,A0t} = 3'b000;
    #5 {ENt,A1t,A0t} = 3'b001;
    #5 {ENt,A1t,A0t} = 3'b010;
    #5 {ENt,A1t,A0t} = 3'b011;
    #5 {ENt,A1t,A0t} = 3'b100;
    #5 {ENt,A1t,A0t} = 3'b101;
    #5 {ENt,A1t,A0t} = 3'b110;
    #5 {ENt,A1t,A0t} = 3'b111;
    #10 $finish;
end
endmodule
\end{lstlisting}

\subsubsection{XDC 约束文件}
\begin{lstlisting}
set_property PACKAGE_PIN M4 [get_ports EN]
set_property PACKAGE_PIN N4 [get_ports in_A1]
set_property PACKAGE_PIN R1 [get_ports in_A0]
set_property PACKAGE_PIN K2 [get_ports out_Y0]
set_property PACKAGE_PIN J2 [get_ports out_Y1]
set_property PACKAGE_PIN J3 [get_ports out_Y2]
set_property PACKAGE_PIN H4 [get_ports out_Y3]

set_property IOSTANDARD LVCMOS33 [get_ports *]
set_property BITSTREAM.CONFIG.UNUSEDPIN PULLNONE [current_design]
\end{lstlisting}

\subsubsection{实验现象}
$\texttt{EN}=0$ 时四路 LED 全灭；$\texttt{EN}=1$ 时，根据 $A_1A_0$ 输入，对应编号 LED 单独点亮。

\newpage
\subsection{任务3：带使能端 3-8 译码器（必做）}
\subsubsection{引脚分配}
使能 \textbf{EN(SW0)}、输入 $A_2$(SW1)、$A_1$(SW2)、$A_0$(SW3)；
输出 $\boldsymbol{y}[7:0]$ 依次绑定：K2,J2,J3,H4,J4,G3,G4,F6。

\subsubsection{真值表}
\begin{table}[htbp]
\centering
\begin{tabular}{ccccccccc}
\toprule
$\text{EN}$ & $A_2$ & $A_1$ & $A_0$ & $y_7$ & $y_6$ & $y_5$ & $y_4$ & $y_3\sim y_0$ \\
\midrule
0  & $x$  & $x$  & $x$  & 0  & 0  & 0  & 0  & 0000 \\
1  & 0  & 0  & 0  & 0  & 0  & 0  & 0  & 0001 \\
1  & 0  & 0  & 1  & 0  & 0  & 0  & 0  & 0010 \\
1  & 0  & 1  & 0  & 0  & 0  & 0  & 0  & 0100 \\
1  & 0  & 1  & 1  & 0  & 0  & 0  & 0  & 1000 \\
1  & 1  & 0  & 0  & 0  & 0  & 0  & 1  & 0000 \\
1  & 1  & 0  & 1  & 0  & 0  & 1  & 0  & 0000 \\
1  & 1  & 1  & 0  & 0  & 1  & 0  & 0  & 0000 \\
1  & 1  & 1  & 1  & 1  & 0  & 0  & 0  & 0000 \\
\bottomrule
\end{tabular}
\end{table}

\subsubsection{Verilog 设计代码（行为级）}
\begin{lstlisting}
module dec3_8(
    input        EN,
    input        A2, A1, A0,
    output reg [7:0] y
);

always @(*) begin
    if(!EN)
        y = 8'b00000000;
    else begin
        case({A2,A1,A0})
            3'b000: y = 8'b00000001;
            3'b001: y = 8'b00000010;
            3'b010: y = 8'b00000100;
            3'b011: y = 8'b00001000;
            3'b100: y = 8'b00010000;
            3'b101: y = 8'b00100000;
            3'b110: y = 8'b01000000;
            3'b111: y = 8'b10000000;
            default:y = 8'b00000000;
        endcase
    end
end
endmodule
\end{lstlisting}

\subsubsection{XDC 约束文件}
\begin{lstlisting}
# 输入引脚
set_property PACKAGE_PIN R2 [get_ports EN]
set_property PACKAGE_PIN M4 [get_ports A2]
set_property PACKAGE_PIN N4 [get_ports A1]
set_property PACKAGE_PIN R1 [get_ports A0]
# 输出引脚 y[7:0]
set_property PACKAGE_PIN F6 [get_ports y[7]]
set_property PACKAGE_PIN G4 [get_ports y[6]]
set_property PACKAGE_PIN G3 [get_ports y[5]]
set_property PACKAGE_PIN J4 [get_ports y[4]]
set_property PACKAGE_PIN H4 [get_ports y[3]]
set_property PACKAGE_PIN J3 [get_ports y[2]]
set_property PACKAGE_PIN J2 [get_ports y[1]]
set_property PACKAGE_PIN K2 [get_ports y[0]]

set_property IOSTANDARD LVCMOS33 [get_ports *]
set_property BITSTREAM.CONFIG.UNUSEDPIN PULLNONE [current_design]
\end{lstlisting}

\subsubsection{实验现象}
$\texttt{EN}=0$ 八路 LED 全灭；$\texttt{EN}=1$ 时，3位输入二进制数对应位置 LED 单独点亮。

\newpage
\subsection{任务4：1位全加器（必做）}
\subsubsection{引脚分配}
输入：加数 $A$(SW0)、被加数 $B$(SW1)、低位进位 $\text{CI}$(SW2)；
输出：和 $S$(K2)、进位 $\text{CO}$(J2)。

\subsubsection{逻辑公式}
\[
S = A \oplus B \oplus \text{CI},\quad \text{CO} = AB + A\cdot \text{CI} + B\cdot \text{CI}
\]

\subsubsection{真值表}
\begin{table}[htbp]
\centering
\begin{tabular}{ccccc}
\toprule
$\text{CI}$ & $A$ & $B$ & $S$ & $\text{CO}$ \\
\midrule
0  & 0 & 0 & 0 & 0 \\
0  & 0 & 1 & 1 & 0 \\
0  & 1 & 0 & 1 & 0 \\
0  & 1 & 1 & 0 & 1 \\
1  & 0 & 0 & 1 & 0 \\
1  & 0 & 1 & 0 & 1 \\
1  & 1 & 0 & 0 & 1 \\
1  & 1 & 1 & 1 & 1 \\
\bottomrule
\end{tabular}
\end{table}

\subsubsection{Verilog 设计代码（数据流）}
\begin{lstlisting}
module full_adder(
    input  A, B, CI,
    output S, CO
);
assign S  = A ^ B ^ CI;
assign CO = (A & B) | (A & CI) | (B & CI);
endmodule
\end{lstlisting}

\subsubsection{XDC 约束文件}
\begin{lstlisting}
set_property PACKAGE_PIN R2 [get_ports A]
set_property PACKAGE_PIN M4 [get_ports B]
set_property PACKAGE_PIN N4 [get_ports CI]
set_property PACKAGE_PIN K2 [get_ports S]
set_property PACKAGE_PIN J2 [get_ports CO]

set_property IOSTANDARD LVCMOS33 [get_ports *]
set_property BITSTREAM.CONFIG.UNUSEDPIN PULLNONE [current_design]
\end{lstlisting}

\subsubsection{实验现象}
拨动三路拨码开关组合，两个 LED 分别对应加法和、进位输出，与真值表一致。

\newpage
\subsection{选做任务1：优先报警电路}
\subsubsection{功能说明}
输入满足 $I_3>I_2>I_1>I_0$（优先级从高到低），输出 $Y_0,Y_1,Y_2$ 驱动LED；任意输入有效时 $Y_0$ 必亮。
\subsubsection{Verilog 代码}
\begin{lstlisting}
module alarm_priority(
    input  I0, I1, I2, I3,
    output reg Y0, Y1, Y2
);
always @(*) begin
    Y0 = 1'b1;
    if(I3) begin
        Y1 = 1'b1; Y2 = 1'b1;
    end
    else if(I2) begin
        Y1 = 1'b1; Y2 = 1'b0;
    end
    else if(I1) begin
        Y1 = 1'b0; Y2 = 1'b1;
    end
    else if(I0) begin
        Y1 = 1'b0; Y2 = 1'b0;
    end
    else begin
        Y0 = 1'b0; Y1 = 1'b0; Y2 = 1'b0;
    end
end
endmodule
\end{lstlisting}

\subsection{选做任务2：4位减法+数码管显示电路}
被减数：$\texttt{SW7}\sim\texttt{SW4}$，减数：$\texttt{SW3}\sim\texttt{SW0}$，结果由共阴极数码管显示十进制数值。

\section{实验核查清单（Checklist）}
\begin{table}[htbp]
\centering
\begin{tabular}{|c|l|c|}
\hline
序号 & 检查项目 & 完成标记 $\boldsymbol{\surd}$ \\
\hline
1 & 工程路径无中文、空格，芯片型号 \verb|xc7a35t-1csg324-1| 配置正确 & \\
\hline
2 & Verilog 源码端口、变量无拼写、语法错误 & \\
\hline
3 & 仿真文件编写完整，仿真波形与真值表完全吻合 & \\
\hline
4 & XDC 约束端口与源码一致，物理引脚匹配 EGO1 手册 & \\
\hline
5 & 逻辑综合 Synthesis 无 Error 报错 & \\
\hline
6 & 布局布线 Implementation 无 Error 报错 & \\
\hline
7 & 成功生成 \texttt{*.bit} 比特流下载文件 & \\
\hline
8 & 开发板 JP2 跳线接 JTAG，USB 接线正常、上电正常 & \\
\hline
9 & Vivado 硬件管理器自动识别 FPGA 芯片 & \\
\hline
10 & 程序下载过程无报错，下载完成 & \\
\hline
11 & 板级实测，LED/数码管逻辑与理论完全一致 & \\
\hline
12 & （可选）Flash 固化成功，切换 QSPI 后断电重启功能正常 & \\
\hline
\end{tabular}
\end{table}

\end{document}